\documentclass[11pt,a4paper]{article}

\usepackage[utf8]{inputenc}
\usepackage[T1]{fontenc}
\usepackage[english]{babel}
\usepackage{lmodern}
\usepackage[a4paper,top=2.5cm,bottom=2.5cm,left=2.8cm,right=2.8cm]{geometry}
\usepackage{booktabs}
\usepackage{array}
\usepackage{tabularx}
\usepackage{xcolor}
\usepackage{graphicx}
\usepackage{enumitem}
\usepackage{fancyhdr}
\usepackage{titlesec}
\usepackage[most]{tcolorbox}
\usepackage[hidelinks,breaklinks]{hyperref}
\usepackage{xurl}
\usepackage{microtype}
\usepackage[english=american]{csquotes}

\definecolor{srcgrey}{gray}{0.42}
\definecolor{rule}{HTML}{2F4858}
\definecolor{tint}{HTML}{EEF3F6}

\hypersetup{colorlinks=true, linkcolor=rule, urlcolor=rule, citecolor=rule,
  pdftitle={Complete Food and Home Cooking},
  pdfauthor={Frederik Berg}}

\newcommand{\co}{CO\textsubscript{2}e}

\titleformat{\section}{\large\bfseries\color{rule}}{\thesection}{0.7em}{}
\titleformat{\subsection}{\normalsize\bfseries}{\thesubsection}{0.6em}{}
\titlespacing*{\section}{0pt}{1.7em}{0.7em}

\setlist[itemize]{topsep=3pt,itemsep=3pt,parsep=0pt,leftmargin=1.2em}
\setlength{\parskip}{0.5em}
\setlength{\parindent}{0pt}

\pagestyle{fancy}
\fancyhf{}
\fancyhead[L]{\footnotesize\color{srcgrey}Complete food and home cooking}
\fancyhead[R]{\footnotesize\color{srcgrey}\thepage}
\renewcommand{\headrule}{\color{rule}\hrule height0.6pt width\headwidth}
\setlength{\headheight}{14pt}

\newcolumntype{R}{>{\raggedleft\arraybackslash}X}
\newcolumntype{L}{>{\raggedright\arraybackslash}X}

\newtcolorbox{keybox}{
  colback=tint, colframe=rule, boxrule=0.6pt, arc=2pt,
  left=9pt, right=9pt, top=8pt, bottom=8pt}

\begin{document}

\thispagestyle{empty}

\begin{center}
{\LARGE\bfseries Complete Food and Home Cooking}\\[0.5em]
{\large Environmental and health evidence for two ways of eating a vegan diet}\\[1.3em]
{\normalsize Frederik Berg}\\[0.3em]
{\footnotesize\color{srcgrey}July 2026}\\[0.5em]
{\footnotesize\color{srcgrey}Written with AI research and drafting assistance;
 see the note on sources and method at the end}
\end{center}

\vspace{0.8em}
\noindent\textcolor{rule}{\rule{\textwidth}{0.8pt}}
\vspace{0.6em}

\section*{Summary}

A category of food exists now that did not exist fifteen years ago: powders
formulated to be a complete diet on their own. Huel, Jimmy Joy, Mana, Soylent.
Mixed with water they supply calories, protein, fat, fibre and the full set of
vitamins and minerals in the amounts health authorities recommend, and most of the
European brands are vegan. Part of the appeal is environmental — nothing spoils,
portions are exact, no cooking is involved.

Whether that beats simply cooking plant food at home has never been tested in the
published literature. It comes out close enough that environmental impact is not a
good reason to pick either one. Both land somewhere between 1 and 3\,kg of \co{}
per 2000\,kcal, roughly one day of food for one adult. For scale, the same measure
for a meat-heavy diet is about 10\,kg \cite{scarborough}. Whether someone eats
plants at all is a factor-of-four decision; which \emph{form} the plants take is a
factor-of-three decision at most, and much of that three is measurement rather than
substance.

Four findings sit behind that.

\textbf{There is almost nothing to check.} No independent assessment of any complete
food has ever been published in a peer-reviewed journal. One manufacturer publishes
an absolute number; the rest publish either nothing or comparisons with no figures
underneath them.

\textbf{The footprint is created by refining the ingredients, not by making the
product.} Whole soybeans carry 1.22\,kg \co{} per kilogram. Turned into protein
isolate they carry 7.23 — about six times as much, with 84 per cent of it coming
from the factory rather than the field \cite{cc-spi}. Mixing the finished powders
together, the step people picture when they imagine industrial food production,
accounts for 0.2 per cent of the total. The assembly really is nearly free; the
ingredients arrive expensive.

\textbf{The vitamin blend is not the rounding error it is assumed to be.} A generic
premix carries 15\,kg \co{} per kilogram — more, per kilogram, than beef. At 2 per
cent of the product by weight it still contributes roughly a sixth of the ingredient
footprint, more than the oil and the rice flour combined \cite{cc-premix}.

\textbf{On health, the evidence does not exist yet.} Nobody has followed anyone
living mainly on one of these products for longer than four weeks. The usual
objection — that they are ultra-processed — does not transfer cleanly, because the
health risks attached to that category concentrate in processed meat and sweetened
drinks rather than in fortified plant products. Meanwhile the diet these products
compete with has a documented weakness of its own: vegans who do not supplement
carefully show markedly higher fracture rates.

What follows from all of this is that the format of the food is the wrong lever.
Four things move the number more than the choice between a shake and a stew: how the
cooking is powered, how much food gets thrown away, how the shopping gets home, and
whether the vegetables were grown in a heated greenhouse in winter.

\vspace{0.4em}
\noindent\textcolor{rule}{\rule{\textwidth}{0.8pt}}

\section{The two things being compared}

A \textbf{complete food} is a powder formulated to supply everything a person needs.
It keeps for a year at room temperature, requires no cooking, and is normally bought
online in bulk. The category began with Soylent in 2013 and now includes Huel in the
UK, Jimmy Joy in the Netherlands, Mana in the Czech Republic, Saturo, YFood and
others. The products examined here are all vegan.

The comparison is against a \textbf{regional, seasonal, whole-food vegan diet}:
staples and produce bought in Germany, weighted toward what is in season or in
domestic cold storage, and cooked at home.

Holding both sides vegan is what makes the question interesting. Removing animal
products is well established as the dominant lever on dietary footprint, so fixing
it leaves a narrower question: given that someone eats plants, does the form those
plants arrive in change anything?

Manufacturers have an obvious commercial interest in the answer, so independent
databases are preferred wherever both exist, and every manufacturer figure is
labelled as one.

\section{Why the numbers are hard to compare}

Food footprints come from \textbf{life-cycle assessment} — an accounting method that
traces every input, from fertiliser and diesel through factory electricity to
packaging and transport, and adds up the emissions attributable to one unit of
product. Two choices decide what such an assessment says, and both are easy to get
wrong.

The first is the \textbf{system boundary}: where the accounting starts and stops.
Cradle-to-gate stops when the product leaves the factory. Cradle-to-shelf continues
to the supermarket. Cradle-to-grave includes cooking, refrigeration and waste at
home. A product measured cradle-to-gate always looks better than the same product
measured cradle-to-grave, so comparing across boundaries produces nonsense.

The second is the \textbf{functional unit}: what the footprint is expressed per. This
one decides the answer here. Complete food is a dry powder at about 417\,kcal per
100\,g. Fresh vegetables are 85 to 95 per cent water. Compare them per kilogram and
the powder looks roughly ten times better — entirely because of the water. A kilogram
of powder is four days of food; a kilogram of tomatoes is a side salad. Everything in
this paper is therefore expressed \textbf{per 2000\,kcal}, and any manufacturer
figure quoted per kilogram should be treated as uninformative until converted.

Three sources carry the quantitative work. \textbf{Agribalyse 3.1}, from the French
environment agency, covers whole foods from cradle to shelf and breaks each figure
down by stage — agriculture, processing, packaging, transport, retail. It is the only
public database listing seasonal and out-of-season produce separately, which is what
makes the seasonality question answerable at all \cite{agribalyse}.
\textbf{CarbonCloud ClimateHub} covers the industrial ingredients Agribalyse does not
— protein isolates, oil powders, vitamin premixes — at factory gate, each split into
agriculture, processing, transport and packaging \cite{climatehub}. And
\textbf{Scarborough and colleagues}, in \emph{Nature Food}, provide the measured
reference point: the actual diets of 55,504 people, each food linked to a review of
570 life-cycle assessments covering 38,000 farms in 119 countries
\cite{scarborough}. That is what people really eat, not a modelled ideal.

Household waste, cooking energy and the trip to the shop are added as separate,
visible lines rather than folded into the food figure, so any comparison here can be
re-run at a different boundary.

\section{What the manufacturers publish}

Very little, and almost none of it independently checkable.

\begin{table}[h]
\centering
\footnotesize
\begin{tabularx}{\textwidth}{@{}lLL@{}}
\toprule
\textbf{Brand} & \textbf{Published product footprint} & \textbf{Verification} \\
\midrule
Huel & 0.6\,kg \co{} per 400\,kcal serving, 83\,\% of it from ingredients & Company data, assured to ISO 14064-3:2019; no methodology published \\
\addlinespace
Jimmy Joy & None. Reports 103\,t \co{} offset in total & Offsetting, not measurement \\
\addlinespace
Mana & None found & --- \\
\addlinespace
Soylent & Comparative ratios only, no absolute figure & Company-commissioned, unpublished \\
\addlinespace
Saturo, YFood, Feed., Queal & None found & --- \\
\bottomrule
\end{tabularx}
\end{table}

Huel is the outlier and deserves credit for it. The company measures its emissions
annually at both company and product level and has the result independently assured
\cite{huel-sustain,huel-nutrition}. Its stated design target is below 0.5\,kg \co{}
per meal, and its help centre describes the measurement as covering ingredients,
manufacturing, packaging and delivery \cite{huel-help}.

That assurance is nonetheless weaker than it sounds. The standard used verifies that
an emissions inventory was compiled correctly. It does not require the company to
disclose which background database it used, how emissions were divided between
products, or where the system boundary was drawn. Without those three things an
outsider cannot reproduce the figure or compare it fairly against a database value.

Jimmy Joy illustrates the more common pattern. Its sustainability page describes
portion control, a twelve-month shelf life, partly recycled packaging and a
tree-planting partnership, and states that the company has offset 103 tonnes of \co{}
in total — which the page itself equates to taking 22 cars off the road for a year
\cite{jimmyjoy}. None of that is a footprint. Offsetting is a payment made after the
emissions have happened, and a cumulative 103 tonnes across the life of a company is
far too small to cover a supply chain that grows and processes ingredients. There is
no measurement there to examine.

The absence is the single most important methodological fact in this area:
\textbf{no independent, peer-reviewed assessment of any complete-food product
exists.} Every quantitative statement about these products, including the ones
below, is either a manufacturer claim or a reconstruction from ingredient data.

\section{Where the footprint actually comes from}

What determines a formulated food's footprint is how heavily its ingredients were
refined before they reached the mixer. The effect is large and it is measurable.

Refining means separating one component of a crop from the rest. Milling a grain into
flour barely counts — the whole seed is still there, just smaller. Pulling the protein
out of a bean and discarding the rest is heavy refining: it needs water, chemicals to
shift the acidity so the protein separates, and then a great deal of heat to dry the
result back into a powder. That energy shows up directly in the footprint.

\begin{center}
\begin{tcolorbox}[colback=tint,colframe=rule,boxrule=0.6pt,arc=2pt,left=10pt,right=10pt,top=8pt,bottom=8pt,width=0.9\textwidth]
\centering
\small
\textbf{whole soybean} \quad$\longrightarrow$\quad \textbf{soy flour} \quad$\longrightarrow$\quad \textbf{soy protein isolate}\\[0.4em]
1.22 \hspace{5.0em} 0.64 \hspace{6.4em} \textbf{7.23}\\[0.5em]
\footnotesize kg \co{} per kg of ingredient. 84\,\% of the last figure is factory energy, not farming.
\end{tcolorbox}
\end{center}

The same pattern repeats across every pair that has been measured.

\begin{table}[h]
\centering
\footnotesize
\begin{tabular}{@{}lrr@{}}
\toprule
\textbf{Ingredient} & \textbf{kg \co{} per kg} & \textbf{Share from processing} \\
\midrule
Soybean & 1.22 & $\approx$0\,\% \\
Soybean meal & 0.64 & 21\,\% \\
Soy protein isolate (90\,\% protein) & \textbf{7.23} & \textbf{84\,\%} \\
\addlinespace
Rolled oats & 1.20 & 2\,\% \\
Oat flour, whole grain & 1.41 & --- \\
Hydrolysed oat flour & \textbf{4.17} & \textbf{77\,\%} \\
\addlinespace
Sunflower oil & 1.73 & 12\,\% \\
Sunflower oil dried into powder & \textbf{3.13} & \textbf{45\,\%} \\
\addlinespace
Pea protein isolate (83\,\% protein) & 3.62 & 60\,\% \\
Vitamin and mineral premix & \textbf{15.0} & \textbf{100\,\%} \\
\bottomrule
\end{tabular}
\end{table}

\noindent All values are CarbonCloud verified product reports, at factory gate
\cite{cc-spi,cc-oat,cc-sun,cc-pea,cc-premix}.

Grinding is close to free. Extracting a protein or drying an oil onto a carrier is
not, and the extra cost is process energy rather than agriculture — which means
sourcing the crop better does not remove it.

This has a consequence neither the manufacturers nor their critics tend to mention:
\textbf{two products in the same category can differ from each other more than
either differs from home cooking.} Jimmy Joy's Plenny Shake is built mainly from
flours, with soy protein isolate appearing only fifth on the ingredient list
\cite{off-plenny}. Huel Powder is built around pea and rice protein isolates and a
dried oil. By the table above, a flour-based formulation should sit near the bottom
of the plausible range and an isolate-based one near the top. No brand publishes
enough to confirm it.

The vitamin premix deserves separate attention, because the usual assumption about
it is wrong. It is dismissed on the grounds that the quantities are tiny — micrograms
of B12, milligrams of most of the rest. That is true of the vitamins and false of the
premix, which is mostly mineral salts and carrier material and which at 15\,kg \co{}
per kilogram is more carbon-intensive per kilogram than beef is per kilogram of meat
\cite{cc-premix}. At roughly 2 per cent of product mass it still contributes about
0.30\,kg \co{} per kilogram of finished powder. That is a line item, not noise.

\section{Estimating a complete food from its ingredients}

Because no product assessment exists, one has to be built. The method is
straightforward, and it works because food labelling is regulated: the ingredient
list must be in descending order by weight and the nutrition panel must be accurate.
Between the two, find the set of ingredient quantities that reproduces the panel
while respecting the ordering, then apply the emission factors above.

Plenny Shake declares, per 100\,g: 417\,kcal, 20.8\,g protein, 15.6\,g fat, 43.8\,g
carbohydrate and 8.6\,g fibre, with the ingredient order oat flour, soy flour,
sunflower oil, rice flour, soy protein isolate, ground golden flaxseed, inulin,
vitamin and mineral mix \cite{off-plenny}. The fitted composition reproduces protein,
fat and fibre to within 2 per cent and carbohydrate to within 9.

\begin{table}[h]
\centering
\footnotesize
\begin{tabular}{@{}lrrrr@{}}
\toprule
\textbf{Component} & \textbf{g/100\,g} & \textbf{kg \co{}/kg} & \textbf{Contribution} & \textbf{Share} \\
\midrule
Oat flour & 52.9 & 1.41 & 0.747 & 40.0\,\% \\
Soy flour & 30.4 & 1.40 & 0.426 & 22.8\,\% \\
Vitamin and mineral premix & 2.0 & 15.0 & 0.300 & 16.1\,\% \\
Soy protein isolate & 2.8 & 7.23 & 0.202 & 10.8\,\% \\
Sunflower oil & 5.2 & 1.73 & 0.090 & 4.8\,\% \\
Rice flour & 4.3 & 1.90 & 0.081 & 4.3\,\% \\
Ground flaxseed & 0.8 & 2.40 & 0.020 & 1.1\,\% \\
\midrule
\textbf{Ingredients} & & & \textbf{1.87} & 100\,\% \\
Dry blending & & & 0.005 & \\
Multilayer pouch & & & 0.12 & \\
Parcel delivery, bulk order & & & 0.15 & \\
\midrule
\textbf{Total} & & & \textbf{2.14\,kg \co{}/kg} & \\
\textbf{Per 2000\,kcal} & & & \textbf{1.03\,kg \co{}} & \\
\bottomrule
\end{tabular}
\end{table}

The central estimate is 1.03\,kg \co{} per 2000\,kcal, with a range of 0.86 to 1.31
across the plausible span of each input.

The line worth pausing on is dry blending at 0.005 — two tenths of one per cent of
the total. That step is the factory: powders going into a mixer and coming out
combined. It is the part people picture when they imagine industrial food production,
and it is almost free. Everything that mattered happened before the ingredients
arrived.

\subsection{A gap that needs explaining}

This estimate is about a third of Huel's self-reported figure, which works out to
3.0\,kg \co{} per 2000\,kcal. The discrepancy is larger than the difference the whole
exercise is trying to measure, so it cannot be left alone.

Part of it is the recipe. Re-running the same model on a Huel-like formulation — oats,
pea protein isolate, flaxseed, rice protein, dried oil — gives roughly 1.4 to 1.5\,kg
\co{} per 2000\,kcal, because isolates and oil powders replace flours. That closes
perhaps half the gap.

The most plausible explanation for the rest is that the figure is not measuring what
it appears to measure. A company can count its emissions two ways: it can assess the
product, or it can total up the whole business — offices, staff travel, marketing,
servers, equipment — and divide by meals sold. The second number is always bigger.
Huel reports carbon intensity per pound of revenue alongside the per-meal figure
\cite{huel-sustain}, and revenue intensity is a whole-business measure rather than a
product one.

If that reading is right, the manufacturer's number is conservative rather than
flattering — which is unusual enough to be worth stating, because the reflexive
assumption about company sustainability figures runs the other way. It remains a
hypothesis; no published methodology would settle it. Both figures are carried into
the comparison below.

\section{The home-cooked side}

Built the same way, from the same databases, so that the two sides can be compared
without adjusting for method. Two days of food were constructed at 2000\,kcal each:
one optimised around German seasonal and stored produce in winter, one reflecting an
ordinary vegan shop including imported and greenhouse items.

\begin{table}[h]
\centering
\footnotesize
\begin{tabular}{@{}lrr@{}}
\toprule
& \textbf{A: optimised} & \textbf{B: ordinary} \\
& \textbf{regional and seasonal} & \textbf{imported + greenhouse} \\
\midrule
Food, cradle-to-shelf & 1.14 & 1.47 \\
Food thrown away at home & +0.14 & +0.22 \\
Cooking and refrigeration & +0.57 & +0.57 \\
Trip to the shop & +0.15 & +0.30 \\
\midrule
\textbf{As eaten, per 2000\,kcal} & \textbf{1.99} & \textbf{2.55} \\
\bottomrule
\end{tabular}
\end{table}

Three results here matter more than the totals.

\textbf{Seasonality is a large effect, and it is about heating, not distance.}
Agribalyse prices a seasonal tomato at 0.626\,kg \co{} per kilogram and an
out-of-season one at 2.114 — a factor of 3.4 \cite{agribalyse}. The stage breakdown
shows where that comes from: the agricultural stage rises from 0.222 to 1.670 while
transport barely moves, from 0.298 to 0.338. The penalty is the heated greenhouse. In
basket B a single 200\,g portion of out-of-season tomato contributes 0.43\,kg \co{},
which is the largest single line in the day and more than all the grain in it put
together. Buying that tomato in summer instead of February takes about a fifth off
the footprint of the whole day's food. Nothing else on the shopping list comes close.

\textbf{Transport is a large share of a small number, which is why food miles
mislead.} For apples, transport is 64 per cent of the cradle-to-shelf footprint — but
that footprint is 0.408\,kg \co{} per kilogram in total. For dried lentils it is 42
per cent of 0.582. The share is high precisely because growing plants is cheap; there
is little else in the number for transport to compete with, so optimising it yields
very little in absolute terms.

This reconciles two studies usually presented as contradicting each other. One found
transport at 11 per cent of food-related emissions, with final delivery from producer
to retailer at 4 \cite{weber}. A later one put global food transport near 19 per cent
\cite{li}. The higher figure counts all transport anywhere in the chain — fertiliser
moving to farms, feed to livestock, machinery to dealers — rather than only finished
food moving to a shop. For someone standing in front of two apples, the lower figure
is the relevant one.

\textbf{Cooking energy is the largest non-food line, and most published comparisons
leave it out.} At 0.57\,kg \co{} per day it is half the footprint of the food itself
in basket A. It is also the one place where a powder mixed with cold tap water has an
unambiguous structural advantage, along with waste: a sealed pouch with measured
scoops does not go off at the back of a fridge.

For calibration against reality, the measured average UK vegan diet is 2.47\,kg \co{}
per 2000\,kcal, with a 95 per cent uncertainty interval of 2.09 to 3.36
\cite{scarborough}. That figure stops at the retail shelf, so it belongs beside the
1.14 and 1.47 rows rather than beside the totals, and it sits above both modelled
baskets — which is expected, because real diets include plant milks, meat
substitutes, coffee, chocolate, snacks and alcohol that a staples-only model leaves
out. Germany's environment agency independently arrives at about 2.85\,kg per day for
a vegan diet with partly regional, seasonal and organic sourcing \cite{uba-food}.

\section{Side by side}

\begin{table}[h]
\centering
\footnotesize
\begin{tabularx}{\textwidth}{@{}LrL@{}}
\toprule
\textbf{Option} & \textbf{kg \co{} per 2000\,kcal} & \textbf{Basis} \\
\midrule
Complete food, built up from ingredients & 1.03 \; (0.86--1.31) & Reconstructed from verified ingredient factors \\
Optimised regional and seasonal winter day & 1.99 & Agribalyse plus German grid \\
Ordinary vegan day & 2.55 & Agribalyse plus German grid \\
Measured average UK vegan diet & 2.47 \; (2.09--3.36) & Scarborough 2023, stops at retail \\
Complete food, manufacturer figure & 3.00 & Company figure, methodology unpublished \\
\addlinespace
\emph{For scale: measured UK high-meat diet} & \emph{10.24 \; (7.04--15.95)} & \emph{Scarborough 2023, stops at retail} \\
\bottomrule
\end{tabularx}
\end{table}

Everything in the plant-based block lands between 1.0 and 3.0. The meat-heavy
reference sits at 10.24. Moving from a meat-heavy diet to any vegan pattern is a
factor of four; every choice available \emph{within} the vegan block is a factor of
three at most, and a good part of that three is measurement rather than substance.

Two rows are not strictly comparable and are marked as such. The measured figures stop
at the retail shelf while the two baskets and the reconstruction run through to
consumption. Comparing like with like, the food-only numbers are 1.14 for the
optimised basket, 1.47 for the ordinary one, 2.47 for the measured average and about
0.90 for the complete food. On that basis the complete food is ahead — but by a margin
comfortably inside the uncertainty of a reconstruction that had to infer its own
recipe.

\section{Impacts other than carbon}

Carbon is the only indicator any complete-food manufacturer reports. None publishes
land use, water scarcity, nutrient pollution or biodiversity impact. That is a gap,
not a clean bill of health.

What can be said from ingredient data is directional. The formulation is dominated by
oats and soy, which are land-hungry per kilogram compared with vegetables but
efficient per calorie. Agribalyse gives oat flakes 214\,m\textsuperscript{2}$\cdot$year
of land per kilogram against carrots at 9.54 — but oats deliver about 380\,kcal per
100\,g and carrots about 40, so per calorie the ranking reverses \cite{agribalyse}.
This is the functional-unit trap from Section~2 running in the other direction, and a
reminder that a single number without its unit is not evidence of anything.

Two effects favour the complete food structurally and are not speculative. Household
waste for a shelf-stable powder in pre-measured servings is effectively zero, against
fresh fruit and vegetables making up around a third of avoidable household food waste
in Germany. And packaging per calorie is low, because the product is dehydrated —
although multilayer laminate pouches are in practice not mechanically recycled in
Germany and go to energy recovery, a circularity failure the carbon number does not
capture.

One effect runs the other way. Ready-to-drink versions ship roughly five times the
mass per calorie, because they ship water, and they need rigid bottles rather than
flexible pouches. The choice between powder and ready-to-drink within the category
plausibly matters more than the choice between the category and cooking.

\section{Health}

The environmental question comes out close. The health question is genuinely
unresolved, and the reason is an absence of evidence rather than a balance of it.

\subsection{What is not known}

No trial or cohort study has followed healthy, weight-stable adults getting most of
their calories from a complete food for longer than four weeks. The longest published
study is a single-arm, industry-funded pilot with twenty people who finished it.

The large body of research on meal replacements does not fill that gap. Those trials
use meal replacements for deliberate weight loss in people with obesity or type 2
diabetes — a different intervention, in a different population, at a different calorie
level.

Specific unknowns worth naming rather than glossing over: long-term effects on teeth
and chewing muscles from a mainly liquid diet, appetite regulation over months, what
low fibre variety does to the gut microbiome, and how well minerals are actually
absorbed from an oat and soy matrix, which is high in phytate — a compound in seeds and
grains that binds iron and zinc and reduces uptake. That last question has never been
measured in these products.

\subsection{The ultra-processed food question}

Complete foods fall squarely into the category most nutrition research now calls
ultra-processed. That category has an unfavourable reputation, and at first glance the
evidence supports it: a 2024 review pooling 45 analyses found higher intake associated
with a 21 per cent higher risk of dying from any cause, a 50 per cent higher risk of
cardiovascular death, and raised risks of type 2 diabetes and depression
\cite{bmj-upf}. Only two of those estimates reached moderate certainty on the standard
grading scale; most remained low or very low, which is normal for observational
nutrition research.

The problem is what the category contains. Processed meat, soft drinks, confectionery
and a fortified vegan shake all sit in the same bucket, and a large European study
shows they do not behave alike. Across 266,666 participants followed for a median of
11.2 years, higher intake did raise the risk of cancer and cardiometabolic disease —
but the effect concentrated in animal-based products and sweetened drinks, while
ultra-processed breads and cereals and plant-based alternatives showed no detectable
association at all \cite{epic-upf}. Complete foods belong to the second group by
composition.

The strongest evidence on the other side is the only randomised trial that isolates
processing itself. Twenty weight-stable adults lived in a research ward and ate an
ultra-processed diet for two weeks and an unprocessed one for two weeks, in random
order. The two menus were matched for calories offered, energy density, protein, fat,
carbohydrate, sugar, salt and fibre, so the only difference was processing. On the
ultra-processed diet people ate 508\,kcal a day more without being asked to and gained
0.9\,kg; on the unprocessed diet they lost 0.9\,kg \cite{hall}.

That is causal evidence that processing drives overeating regardless of nutrient
content. How far it transfers is arguable in both directions — the menu in that trial
was conventional packaged food, not a fortified formulation, but the mechanism
identified, eating faster with less satiety per calorie, applies at least as strongly
to something drinkable. Nobody has run the experiment with a complete food.

\subsection{The risk on the other side}

It is easy to treat home-cooked vegan food as the safe default. The data does not
entirely support that.

In a cohort of at least 54,898 people followed for an average of 17.6 years, vegans
had 2.31 times the hip-fracture rate of meat eaters, with a 95 per cent confidence
interval of 1.66 to 3.22. In absolute terms that is 14.9 additional hip fractures per
1000 people over ten years. Total fracture risk was also raised, at 1.43
\cite{epic-fracture}. The associations weakened but persisted after adjusting for
dietary calcium and protein, which points at intake adequacy rather than at anything
intrinsic to plant food.

Micronutrient status tells the same story from another angle. Vitamin B12 is the
familiar case; iodine, selenium, zinc and long-chain omega-3 fats are the
less-discussed ones, and in European populations — where soils are low in selenium and
salt is not routinely iodised — they are not rare deficiencies.

This is the honest case in favour of complete food, and it is stronger than the
environmental one. A product supplying calcium, vitamin D, B12, iodine, selenium, zinc
and algal omega-3 at reference levels removes, by construction, the exact failure the
fracture data points at. Whether it does so in practice has never been tested, and
\enquote{should eliminate} is not \enquote{eliminates}.

\subsection{Contaminants}

Plant proteins concentrate heavy metals from soil more than animal proteins do, and
organic sourcing tends to increase rather than reduce cadmium and lead uptake.
Exceedances of California's Proposition 65 thresholds are common in commercial protein
powders, though those thresholds are far stricter than European limits and presence at
those levels is not evidence of harm.

The consideration is real for a product eaten daily in quantity rather than
occasionally, and it is a place where the varied diet has a genuine advantage: eating
many different plants dilutes any single contaminated input.

\section{The strongest objection}

The case against the conclusion runs like this. The comparison flatters complete food,
because the home-cooked baskets carry cooking, shopping and waste while the powder
carries none of them. Those three lines are 0.86 of the 1.99\,kg optimised total —
nearly half. Strip them out and the food itself is 1.14 against about 0.90, close to
parity and well inside the error bars of a reconstruction built on an inferred recipe.

That objection is correct, and it is the conclusion rather than a refutation of it.
The lines are not an artefact; they are real emissions a powder user genuinely does
not produce. But their size relative to the food is the point. What separates these
two options is not the food, it is everything around the food. Someone who cooks on an
induction hob with renewable electricity, shops by bicycle and wastes nothing has
already closed the gap without changing what they eat.

A second objection is that the reconstruction is the weakest number here and is doing
the most work. Also correct. The recipe is inferred rather than disclosed, the emission
factors come from a commercial database rather than peer review, and the only
manufacturer figure available disagrees with it threefold. It is presented as the
better estimate because its inputs are individually checkable and the manufacturer's
are not. A published product assessment would supersede it immediately.

\section{Established, likely, unknown}

\textbf{Established.} Refining, not manufacturing, drives the footprint of formulated
foods: soy protein isolate costs about six times its parent bean, 84 per cent of that
from processing. Seasonality is a factor of 3.4 for greenhouse-sensitive produce,
driven by heating rather than distance. Transport is a large share of a small absolute
number for plant foods. Cooking energy is a major line and is usually omitted. Vegan
diets of any construction come to roughly a quarter of a meat-heavy diet's footprint.
Vegans who do not supplement carefully carry a measured, substantial excess fracture
risk.

\textbf{Likely.} Complete food probably has a modestly lower as-eaten carbon footprint
than home cooking, driven mainly by zero cooking energy and zero waste rather than by
the food itself. The vitamin premix is a real contributor at roughly a sixth of the
ingredient footprint. Flour-based formulations are materially lighter than
isolate-based ones. The health risks attached to ultra-processed food probably do not
transfer wholesale to fortified plant formulations — though this rests on subgroup
analysis rather than direct study.

\textbf{Unknown.} Any peer-reviewed assessment of a complete-food product. The system
boundary behind the one published manufacturer figure. Land, water, nutrient-pollution
and biodiversity impacts of any of them. Health outcomes beyond four weeks. How well
minerals are absorbed from a high-phytate formulated matrix. Whether fortification
removes the fracture risk it appears designed to remove. Long-term effects on the gut
microbiome, appetite and oral health.

\section{Conclusion}

Environmental impact is not a good basis for choosing between these two ways of eating.
Both sit in the 1 to 3\,kg \co{} band per 2000\,kcal, and the levers that actually move
the number — avoiding out-of-season greenhouse produce, not wasting food, not driving to
the shop, and how the cooking is powered — cut across both.

The intuition that centralising production and adding all the vitamins at once must be
efficient is half right. The assembly really is nearly free: mixing the powders is 0.2
per cent of the product's footprint. But the ingredients arriving at that mixer have
been refined, and refining is where the energy went. Central manufacturing does not
escape that cost; it depends on it.

On health the position is not a tie but a blank. There is no long-term evidence in
either direction, and reading that blank as reassurance is as wrong as reading it as
alarm. What is documented is that the whole-food vegan diet has a measurable failure
mode — micronutrient adequacy, with fracture risk as its visible endpoint — that a
fortified product is designed to solve and has never been shown to solve.

Conditionally, then. Where the motivation is environmental, the evidence does not
support switching, and the same effort spent on seasonality, waste and cooking energy
would achieve more. Where the motivation is convenience, cost or reliable micronutrient
coverage, the environmental evidence neither supports nor blocks it — and partial
substitution captures most of the waste and cooking-energy advantage while keeping the
dietary variety whose long-term value is, unlike most of the above, not seriously
disputed.

\section*{Appendix: basket composition}

Both baskets normalised to 2000\,kcal. Emission factors are Agribalyse 3.1
cradle-to-shelf values in kg \co{} per kg; the last column is the resulting per-day
contribution. Waste, cooking and shopping are added on top of these subtotals as shown
in Section~6. Cooking assumes 1.6\,kWh per day at the German grid factor of 0.353\,kg
\co{} per kWh.

\vspace{0.4em}
\noindent
\begin{minipage}[t]{0.485\textwidth}
\footnotesize
\textbf{A: optimised regional and seasonal winter day}\\[0.3em]
\begin{tabular}{@{}lrrr@{}}
\toprule
Item & g & Factor & kg \co{} \\
\midrule
Potatoes & 350 & 0.788 & 0.276 \\
Rye/wheat bread & 180 & 0.760 & 0.137 \\
Oat flakes & 100 & 1.220 & 0.122 \\
Kale & 150 & 0.713 & 0.107 \\
Walnuts & 25 & 4.110 & 0.103 \\
Sauerkraut & 100 & 0.700 & 0.070 \\
Rapeseed oil & 25 & 2.500 & 0.062 \\
Apple, stored & 150 & 0.408 & 0.061 \\
Leek & 100 & 0.611 & 0.061 \\
Carrots & 150 & 0.396 & 0.059 \\
Beetroot & 100 & 0.408 & 0.041 \\
Lentils, dry & 70 & 0.581 & 0.041 \\
\midrule
\textbf{Subtotal} & & & \textbf{1.14} \\
\bottomrule
\end{tabular}
\end{minipage}
\hfill
\begin{minipage}[t]{0.485\textwidth}
\footnotesize
\textbf{B: ordinary day, imported and greenhouse}\\[0.3em]
\begin{tabular}{@{}lrrr@{}}
\toprule
Item & g & Factor & kg \co{} \\
\midrule
Tomato, off-season & 200 & 2.114 & 0.433 \\
Rice, dry & 80 & 2.015 & 0.165 \\
Tofu & 150 & 1.000 & 0.154 \\
Banana & 150 & 0.909 & 0.140 \\
Bread & 150 & 0.760 & 0.117 \\
Walnuts & 25 & 4.110 & 0.105 \\
Oat flakes & 80 & 1.220 & 0.100 \\
Cucumber & 150 & 0.512 & 0.079 \\
Strawberry, off-season & 100 & 0.592 & 0.061 \\
Sunflower oil & 25 & 2.358 & 0.060 \\
Chickpeas, dry & 60 & 0.899 & 0.055 \\
\midrule
\textbf{Subtotal} & & & \textbf{1.47} \\
\bottomrule
\end{tabular}
\end{minipage}

\vspace{1.0em}
\noindent\textcolor{rule}{\rule{\textwidth}{0.6pt}}
\vspace{0.3em}

{\footnotesize\color{srcgrey}
\textbf{Data and method.} Agribalyse 3.1 queried through the ADEME data API, July
2026. CarbonCloud ClimateHub product reports and manufacturer pages read directly,
July 2026. Recipe fitting by constrained random search over the simplex, minimising
squared relative error against the declared nutrition panel subject to the declared
ingredient ordering. All effect sizes and confidence intervals were checked against
the primary publications rather than against secondary summaries. Uncertainty ranges
reflect the span of published emission factors for each ingredient, not a formal
propagation. Research and drafting were carried out with assistance from Claude
(Anthropic); every source cited was retrieved and checked against the primary
document, and responsibility for the content rests with the author.}

\begin{thebibliography}{99}
\footnotesize

\bibitem{scarborough} Scarborough, P. et al. \emph{Vegans, vegetarians, fish-eaters and meat-eaters in the UK show discrepant environmental impacts.} Nature Food, 2023. \url{https://www.nature.com/articles/s43016-023-00795-w} (open version: \url{https://pmc.ncbi.nlm.nih.gov/articles/PMC10365988/})

\bibitem{agribalyse} ADEME. \emph{Agribalyse 3.1 -- synthèse.} \url{https://data.ademe.fr/datasets/agribalyse-31-synthese}

\bibitem{climatehub} CarbonCloud. \emph{ClimateHub product reports.} \url{https://apps.carboncloud.com/climatehub}

\bibitem{cc-spi} CarbonCloud ClimateHub. \emph{Soy protein isolate, verified product report.} \url{https://apps.carboncloud.com/climatehub/product-reports/id/96565698940}

\bibitem{cc-oat} CarbonCloud ClimateHub. \emph{Oat flour products, verified product report.} \url{https://apps.carboncloud.com/climatehub/product-reports/id/385367484644}

\bibitem{cc-sun} CarbonCloud ClimateHub. \emph{Sunflower oil products, verified product report.} \url{https://apps.carboncloud.com/climatehub/product-reports/id/8670992232269}

\bibitem{cc-pea} CarbonCloud ClimateHub. \emph{Pea protein isolate, verified product report.} \url{https://apps.carboncloud.com/climatehub/product-reports/id/2282851493418}

\bibitem{cc-premix} CarbonCloud ClimateHub. \emph{Vitamin and mineral premix, verified product report.} \url{https://apps.carboncloud.com/climatehub/product-reports/id/186640392289}

\bibitem{off-plenny} Open Food Facts. \emph{Plenny Shake, Jimmy Joy -- ingredients and nutrition panel.} \url{https://world.openfoodfacts.org/product/8720165350995/plenny-shake-jimmy-joy}

\bibitem{huel-nutrition} Huel. \emph{Sustainability nutrition facts.} \url{https://de.huel.com/pages/huel-sustainability-nutrition-facts}

\bibitem{huel-sustain} Huel. \emph{Sustainability.} \url{https://huel.com/pages/sustainability-page}

\bibitem{huel-help} Huel Help Centre. \emph{What is Huel's carbon footprint per meal?} \url{https://intercom-help.eu/uk-help-centre-huel/en/articles/633648-what-is-huel-s-carbon-footprint-per-meal}

\bibitem{jimmyjoy} Jimmy Joy. \emph{How sustainable is Jimmy Joy?} \url{https://jimmyjoy.com/pages/how-sustainable-is-jimmy-joy}

\bibitem{uba-food} Umweltbundesamt. \emph{Klimafreundliche Ernährung.} \url{https://www.umweltbundesamt.de/umwelttipps-fuer-den-alltag/klimafreundliche-ernaehrung-fleischreduziert}

\bibitem{weber} Weber, C.L., Matthews, H.S. \emph{Food-miles and the relative climate impacts of food choices in the United States.} Environmental Science \& Technology, 2008. \url{https://pubs.acs.org/doi/10.1021/es702969f}

\bibitem{li} Li, M. et al. \emph{Global food-miles account for nearly 20\% of total food-systems emissions.} Nature Food, 2022. \url{https://www.nature.com/articles/s43016-022-00531-w}

\bibitem{bmj-upf} Lane, M.M. et al. \emph{Ultra-processed food exposure and adverse health outcomes: umbrella review of epidemiological meta-analyses.} BMJ, 2024. \url{https://pmc.ncbi.nlm.nih.gov/articles/PMC10899807/}

\bibitem{epic-upf} Cordova, R. et al. \emph{Consumption of ultra-processed foods and risk of multimorbidity of cancer and cardiometabolic diseases.} The Lancet Regional Health -- Europe, 2023. \url{https://pmc.ncbi.nlm.nih.gov/articles/PMC10730313/}

\bibitem{hall} Hall, K.D. et al. \emph{Ultra-processed diets cause excess calorie intake and weight gain: an inpatient randomized controlled trial of ad libitum food intake.} Cell Metabolism, 2019. \url{https://pmc.ncbi.nlm.nih.gov/articles/PMC7946062/}

\bibitem{epic-fracture} Tong, T.Y.N. et al. \emph{Vegetarian and vegan diets and risks of total and site-specific fractures: results from the prospective EPIC-Oxford study.} BMC Medicine, 2020. \url{https://bmcmedicine.biomedcentral.com/articles/10.1186/s12916-020-01815-3}

\end{thebibliography}

\end{document}
